\documentclass[11pt,oneside]{article}

\usepackage[margin=0.9in]{geometry}
\usepackage{booktabs}
\usepackage{tabularx}
\usepackage{array}
\usepackage{microtype}
\usepackage{xcolor}
\usepackage{hyperref}
\usepackage{enumitem}
\usepackage{listings}
\usepackage{fancyhdr}
\usepackage{titlesec}
\usepackage{amsmath}
\usepackage{amssymb}

\definecolor{ink}{HTML}{17212B}
\definecolor{accent}{HTML}{075985}
\definecolor{muted}{HTML}{5B6773}
\definecolor{panel}{HTML}{EFF6F9}
\definecolor{rulegray}{HTML}{CBD5E1}
\hypersetup{
  colorlinks=true,
  linkcolor=accent,
  urlcolor=accent,
  citecolor=accent,
  pdftitle={A Cross-Harness Evaluation of Coding-Agent Deployment Systems},
  pdfauthor={Cross-harness coding-agent study}
}

\setlength{\parindent}{0pt}
\setlength{\parskip}{0.58em}
\setlist{nosep,leftmargin=1.5em}
\renewcommand{\arraystretch}{1.16}
\pagestyle{fancy}
\fancyhf{}
\fancyhead[L]{\small\color{muted}Coding-Agent Deployment Systems}
\fancyhead[R]{\small\color{muted}August 2026}
\fancyfoot[C]{\thepage}
\renewcommand{\headrulewidth}{0.3pt}
\titleformat{\section}{\Large\bfseries\color{ink}}{\thesection}{0.7em}{}
\titleformat{\subsection}{\large\bfseries\color{accent}}{\thesubsection}{0.7em}{}
\sloppy

\lstdefinestyle{config}{
  basicstyle=\ttfamily\small,
  backgroundcolor=\color{panel},
  frame=single,
  rulecolor=\color{rulegray},
  breaklines=true,
  columns=fullflexible,
  showstringspaces=false
}

\newcommand{\resultbox}[1]{%
  \begin{center}
  \fcolorbox{accent}{panel}{\parbox{0.92\textwidth}{#1}}
  \end{center}
}

\begin{document}

\hypersetup{pageanchor=false}
\begin{titlepage}
  \centering
  \vspace*{1.5cm}
  {\Huge\bfseries\color{ink}Coding-Agent\\[0.2em]Deployment Systems\par}
  \vspace{0.7cm}
  {\LARGE A Predeclared Cross-Harness Evaluation\\with a Post-Hoc Luna Solo Supplement\par}
  \vspace{1.2cm}
  {\large A thesis-style technical report and corrective extension\par}
  \vfill
  \begin{tabular}{rl}
    Date: & 15 August 2026 \\
    Host: & Linux x86\_64, 4 logical CPUs, 15.5 GiB RAM \\
    Cursor Agent CLI: & \texttt{2026.08.11-e8db854}, authenticated subscription \\
    OpenCode: & 1.18.18, authenticated subscription \\
    New outcome runs: & 35 predeclared plus 10 post-hoc exploratory \\
    Prior controlled runs: & 40 in OpenCode and Pi \\
  \end{tabular}
  \vfill
  {\color{muted}\small
  The complete new protocol, raw streams, hidden-grader outputs, hashes, and
  analysis are under \texttt{reproducibility/cross-harness/}.
  }
\end{titlepage}
\hypersetup{pageanchor=true}

\begin{abstract}
An earlier experiment on this machine found that a GPT-5.6 Sol/xhigh parent
with three bounded Terra/medium workers completed one separable coding workload
faster than Sol/xhigh alone in OpenCode and Pi. That study omitted a strong
external baseline: a fast single-agent system in Cursor CLI.

This extension predeclared seven deployment arms in a local protocol frozen
before outcomes and ran five randomized complete blocks, producing 35 fresh
outcome runs on the same hidden-graded, three-package workload. The arms were
solo Cursor Grok 4.5 high, Grok 4.6 high,
and Grok 4.6 xhigh; OpenCode Sol/xhigh solo; bounded Sol+Terra/medium; bounded
Sol+Luna/max; and production-like Sol+Luna/max. Cursor's Task tool was excluded,
and raw streams confirmed zero subagent calls. The workload was deliberately
separable, favoring parallel delegation.

The pre-outcome freeze was local rather than an externally archived immutable
preregistration. Current analysis checks artifact and score integrity, but cannot
prove that the source files were unchanged before execution.

Solo Grok 4.5 high was fastest and highest-scoring: 113/120 hidden checks in
109.5 seconds mean agent time. OpenCode Sol/xhigh solo scored 112/120 in 486.8
seconds. Grok 4.5 reduced aggregate agent time by 77.5\%; every block favored
Grok, and the mean paired reduction was 77.0\% with a descriptive 95\%
small-sample interval of 69.2--84.7\%. Bounded Terra scored 106/120 in 405.1
seconds. Bounded and natural Luna scored 110/120 in 949.4 and at least 1190.2
seconds mean; natural Luna timed out once. Grok 4.6 high and xhigh did not
improve aggregate quality and were 4.0 and 5.7 times slower than Grok 4.5.

After the parent analysis, a requested post-hoc supplement randomized Luna/max
solo and Luna/xhigh solo in five new blocks. Luna/xhigh scored 110/120 in 750.9
seconds mean agent time; Luna/max scored 109/120 in 779.4 seconds. Xhigh used
3.7\% less aggregate time but was faster in only two blocks, and its paired
interval crossed zero. Against non-concurrent historical outcomes, Grok used
85.4--86.0\% less aggregate time and Sol used 35.2--37.5\% less, with higher
quality. This corroborates but does not alter the frozen parent selection.

No arm achieved a perfect run. The result is therefore not permission to skip
verification. It is a deployment decision: default to solo Cursor/Grok 4.5 high
for the tested workflow; use OpenCode Sol/xhigh solo if OpenCode is required;
do not replace Sol with Luna based on the post-hoc solo screen; disable Luna/max
subagents; and treat bounded Terra as experimental rather than a default.
\end{abstract}

\tableofcontents
\newpage

\section{Decision}

\resultbox{\textbf{Operating recommendation.} Default to solo Cursor Agent with
Grok 4.5 high. Keep Task/subagent invocation disabled unless a task-specific,
hidden-graded comparison demonstrates an advantage. If OpenCode is required,
use Sol/xhigh solo with \texttt{task} denied. Disable Luna/max subagents. Treat
Luna solo as unsupported by the post-hoc screen. Treat bounded Terra as an
experimental low-risk mode, not a security-sensitive default.}

This practical production command distills the tested default:

\begin{lstlisting}[style=config]
agent -p --force --trust \
  --model cursor-grok-4.5-high \
  --output-format stream-json \
  --exclude-tools task_tool_call \
  "<task>"
\end{lstlisting}

It is not the exact isolated invocation. The harness also used a fresh
\texttt{CURSOR\_CONFIG\_DIR} and \texttt{--data-dir}, explicit workspace,
\texttt{--single-turn}, \texttt{--sandbox disabled}, and flags disabling
project config, indexing, and codebase references. Exact reproduction uses
\texttt{run.mjs}.

The installed CLI accepted \texttt{--exclude-tools task\_tool\_call}, although
that flag was absent from ordinary help output. Cursor documents automatic
subagents, so model selection alone does not prove solo execution. The five
outcome streams each contained one session ID and zero Task calls.

This policy is evidence-calibrated. It is strong for this fixture and observation
window, moderate as a personal default, and not a universal model ranking.
The post-hoc solo supplement also found historical Sol 35.2--37.5\% faster than
Luna/xhigh or Luna/max, with higher aggregate quality.

\section{Why the Earlier Thesis Was Incomplete}

The prior study asked whether subagents improved an OpenCode-centered workflow.
Its within-harness findings were useful:

\begin{center}
\begin{tabularx}{\textwidth}{@{}Xrrrr@{}}
\toprule
Prior harness & Single mean & Delegated mean & Time change & Quality \\
\midrule
OpenCode 1.18.7, five pairs & 19.5m & 11.3m & $-42.1\%$ & 108/120 vs 112/120 corrected \\
Pi 0.80.6, three pairs & 12.2m & 8.0m & $-33.9\%$ & 68/72 vs 66/72 \\
\bottomrule
\end{tabularx}
\end{center}

Those results answer an architecture question conditional on using Sol and the
tested harness. They do not answer whether to use OpenCode and GPT-5.6 at all.
Sol/xhigh was treated as an established default rather than tested against
Cursor/Grok. The previous report disclosed that limitation but still made an
operating recommendation beyond its evidence.

This is not a cosmetic distinction. SWE-agent shows that agent-computer
interface design can materially affect coding performance, while Agentless
demonstrates that simple systems can outperform elaborate autonomous agents.
The intervention relevant to a user is the complete deployment system: harness,
model, tools, reasoning effort, and delegation policy.

\section{Research Questions}

\begin{description}
  \item[RQ1] Which tested solo Cursor/Grok configuration is quality-eligible
  and fastest on the frozen workload?
  \item[RQ2] Does the best solo Cursor arm beat OpenCode Sol/xhigh alone on
  hidden quality and user-perceived completion time?
  \item[RQ3] Does it beat the earlier recommended bounded Sol+Terra system?
  \item[RQ4] Within current OpenCode, do Terra or Luna subagents improve time
  without unacceptable quality loss?
  \item[RQ5] Do the tested higher-effort configurations improve quality enough
  to justify their latency?
  \item[RQ6] What operating policy survives the combined old and new evidence?
\end{description}

After the parent study was complete, the user requested a separate exploratory
question: how do Luna/max and Luna/xhigh perform as solo OpenCode owners? It is
not retroactively designated a preregistered research question.

The estimand is deployable-system performance. This study does not identify a
pure causal effect of concurrency or model family.

Capacity was asymmetric by design: OpenCode solo had 36 parent turns; bounded
arms had 12 parent turns plus up to 36 child turns; natural Luna could continue
until the wall cap. Cursor exposed no equivalent turn cap. These are usable
architectures, not equal-compute treatments.

\section{Preregistered Design}

\subsection{Experimental arms}

\begin{center}
\small
\begin{tabularx}{\textwidth}{@{}XllX@{}}
\toprule
Arm & Parent / solo & Children & Delegation policy \\
\midrule
Cursor Grok 4.5 & Grok 4.5 high & None & Task excluded \\
Cursor Grok 4.6 high & Grok 4.6 high & None & Task excluded \\
Cursor Grok 4.6 xhigh & Grok 4.6 xhigh & None & Task excluded \\
OpenCode solo & Sol/xhigh, 36 turns & None & Task denied \\
Bounded Terra & Sol/xhigh, 12 turns & 3 Terra/medium & One package each; no recursion \\
Bounded Luna & Sol/xhigh, 12 turns & 3 Luna/max & One package each; no recursion \\
Natural Luna & Sol/xhigh & Luna/max & Broad optional recursion \\
\bottomrule
\end{tabularx}
\end{center}

The natural arm reproduced the relevant former production configuration: only
\texttt{general}, \texttt{explore}, and \texttt{scout}, all Luna/max, were
available as subagents.

\subsection{Workload and grading}

Every run received a fresh copy of three independent JavaScript packages:

\begin{itemize}
  \item bounded-concurrency mapping implementation;
  \item session and authentication security repair;
  \item event-pipeline comprehension producing eight exact facts.
\end{itemize}

Each package contributed eight hidden checks, 24 per run. Agents saw public
tests. Graders were not copied into their workspaces, and exact grader paths did
not appear in streams, but OS-level non-access was not enforced. The
comprehension deliverable was fixed before outcomes at
\texttt{comprehension/answer.json}; there was no score-improving path fallback.

The explicit package independence favors subagents. A solo winner therefore
wins despite a workload designed to expose parallelism. The study says little
about coupled refactors.

\subsection{Sampling, randomization, and endpoints}

Five randomized complete blocks were frozen before outcomes. Each contained all
seven arms once. A seeded generator fixed the 35-run order; top-level runs were
sequential. The prompt, fixture hash, arm set, model IDs, five-block sample,
30-minute cap, and decision rule were frozen before outcome inspection.

Primary quality was hidden checks passed. Primary performance was process wall
time from agent launch to exit. Verified wall time added deterministic grading
and produced nearly identical estimates. Checks were clustered within runs and
were not treated as independent samples.

  The predeclared operating rule required no timeout or abnormal exit, no model
mismatch or forbidden final fixture change, aggregate quality within one check
of the best arm, and at least 7/8 security checks in every run. Among eligible
systems, lowest aggregate verified time won.

\subsection{Pilot amendments}

Two excluded pilots found harness defects before outcomes. Cursor rejected
\texttt{--exclude-workspace-context} for this account, so it was removed for all
Cursor arms. The first natural-Luna pilot could see benchmark-only worker
definitions, so that arm was restricted to the user's three production agent
names. Failed pilot artifacts are retained.

\subsection{Post-hoc Luna solo supplement}

On August 15, after the parent result was known, a separate manifest froze
Luna/max solo and Luna/xhigh solo before their own outcomes. Both used OpenCode,
a 36-step cap, denied Task, and received the same fixture, prompt, graders, and
1800-second timeout. Five seeded blocks randomized these two arms and produced
ten sequential outcomes.

Max versus xhigh is randomized within the supplement. Comparisons to August 14
Grok and Sol are non-concurrent descriptive comparisons, not paired effects.
The supplement does not change the parent manifest or operating-rule winner.
Its 43 raw artifacts are hash-pinned; source/config hashes are an explicitly
post-outcome snapshot, not proof of a pre-outcome source freeze.

\section{Main Results}

\subsection{Arm-level outcomes}

\begin{center}
\small
\begin{tabularx}{\textwidth}{@{}Xrrrrr@{}}
\toprule
Arm & Checks & Mean & Median & Range & Timeouts \\
\midrule
\textbf{Grok 4.5 solo} & \textbf{113/120} & \textbf{109.5s} & \textbf{107.0s} & 89--140s & 0 \\
Grok 4.6 high solo & 112/120 & 433.9s & 427.4s & 398--471s & 0 \\
Grok 4.6 xhigh solo & 111/120 & 622.8s & 592.9s & 558--688s & 0 \\
Sol/xhigh solo & 112/120 & 486.8s & 468.3s & 401--574s & 0 \\
Sol+Terra bounded & 106/120 & 405.1s & 402.0s & 230--584s & 0 \\
Sol+Luna bounded & 110/120 & 949.4s & 880.9s & 790--1346s & 0 \\
Sol+Luna natural & 110/120 & $\geq$1190.2s & 1083.9s & 893 to $>1800$s & 1 \\
\bottomrule
\end{tabularx}
\end{center}

Grok 4.5, Grok 4.6 high, and Sol solo passed the fixed gate. Grok 4.5 won
mechanically on time. Grok 4.6 xhigh, Terra, and both Luna arms fell outside the
one-check quality margin; natural Luna also timed out.

\subsection{Grok 4.5 versus OpenCode Sol solo}

\begin{center}
\begin{tabular}{@{}rrrrr@{}}
\toprule
Trial & Grok 4.5 & Sol solo & Time saved & Scores \\
\midrule
1 & 89.3s & 553.4s & 83.9\% & 22 vs 22 \\
2 & 89.5s & 401.0s & 77.7\% & 23 vs 22 \\
3 & 122.1s & 468.3s & 73.9\% & 23 vs 22 \\
4 & 139.8s & 437.6s & 68.0\% & 22 vs 23 \\
5 & 107.0s & 573.6s & 81.3\% & 23 vs 23 \\
\bottomrule
\end{tabular}
\end{center}

Aggregate agent-time reduction was 77.5\%. Mean paired reduction was 77.0\%,
with descriptive 95\% interval 69.2--84.7\%. Mean absolute savings were 377.3
seconds, interval 275.0--479.5 seconds. Every block favored Grok. Aggregate
quality was 113/120 versus 112/120: parity under the fixed margin, not evidence
of general intelligence superiority. No speed-quality tradeoff was observed
under the prespecified margin.

\subsection{All preregistered contrasts}

Block-level agent-time reductions for the candidate named first:

\begin{center}
\small
\begin{tabular}{@{}lrrrrr@{}}
\toprule
Contrast & T1 & T2 & T3 & T4 & T5 \\
\midrule
Grok 4.5 vs Sol & 83.9\% & 77.7\% & 73.9\% & 68.0\% & 81.3\% \\
Grok 4.5 vs Terra & 81.9\% & 77.7\% & 61.5\% & 39.2\% & 81.7\% \\
Terra vs Sol & 11.1\% & $-0.2$\% & 32.2\% & 47.5\% & $-1.8$\% \\
Terra vs Luna & 44.7\% & 52.1\% & 59.9\% & 82.9\% & 33.7\% \\
Sol vs natural Luna & 50.2\% & 62.2\% & 47.6\% & $\geq75.7$\% & 47.1\% \\
\bottomrule
\end{tabular}
\end{center}

\begin{center}
\scriptsize
\begin{tabularx}{\textwidth}{@{}Xrrrccc@{}}
\toprule
Contrast & Aggregate & Mean [95\% interval] & Quality & Faster & Perfect & Timeouts \\
\midrule
Grok 4.5 vs Sol & 77.5\% & 77.0 [69.2,84.7] & 113--112 & 5/5 & 0--0 & 0--0 \\
Grok 4.5 vs Terra & 73.0\% & 68.4 [45.6,91.2] & 113--106 & 5/5 & 0--0 & 0--0 \\
Terra vs Sol & 16.8\% & 17.7 [$-8.9$,44.4] & 106--112 & 3/5 & 0--0 & 0--0 \\
Terra vs Luna & 57.3\% & 54.7 [31.7,77.7] & 106--110 & 5/5 & 0--0 & 0--0 \\
Sol vs natural Luna & $\geq59.1$\% & 56.6 [41.2,71.9] & 112--110 & 5/5 & 0--0 & 0--1 \\
\bottomrule
\end{tabularx}
\end{center}

Natural-Luna trial 4 is right-censored. Its reduction is a lower bound; the
interval using the timeout cap is descriptive, not a completion-time interval.

Mean absolute agent-time savings with 95\% intervals were: Grok 4.5 versus Sol,
377.3 seconds [275.0,479.5]; Grok 4.5 versus Terra, 295.5 [102.1,489.0]; Terra
versus Sol, 81.7 [$-36.8$,200.2]; Terra versus Luna, 544.3 [138.7,949.8]; and
Sol versus natural Luna, 703.4 [233.9,1173.0] using the capped timeout and thus
not an exact completion-time interval.

\subsection{Grok 4.5 versus bounded Terra}

Grok was faster in all five blocks. Aggregate agent-time reduction was 73.0\%;
mean paired reduction was 68.4\%, interval 45.6--91.2\%. Quality was 113/120
versus 106/120. Terra scored only 5/8 on security in two runs and failed the
safety gate. The missing external baseline therefore reverses the old operating
recommendation.

\subsection{Delegation within current OpenCode}

Bounded Terra reduced aggregate agent time by 16.8\% versus Sol solo but was
faster in only three of five blocks. Mean paired reduction was 17.7\%, interval
$-8.9$--44.4\%. Terra lost six checks, 106/120 versus 112/120. Excluding its
transient-test-edit deviation yielded a similar clean-block estimate: 20.1\%
aggregate reduction, interval $-12.6$--57.1\%, and 84/96 versus 90/96 checks.

Terra was 57.3\% faster than bounded Luna but lost four checks. Excluding both
bounded protocol-deviation blocks left three pairs: Terra remained 45.6\%
faster and scored 64/72 versus 66/72.

Sol solo was faster than natural Luna in every block. The capped intention-to-
treat advantage was at least 59.1\%; one natural run is right-censored at 1800
seconds. Across the four normally completed pairs, Sol remained 51.9\% faster
and scored 89/96 versus 88/96.

\subsection{Higher-effort configurations did not dominate}

Grok 4.5 reduced time by 74.8\% versus Grok 4.6 high and 82.4\% versus Grok 4.6
xhigh while scoring 113/120 versus 112/120 and 111/120. This changes model
generation and effort, so it does not isolate a causal effort effect. The
earlier same-model screen also found no gain from Terra/high over Terra/medium,
while Luna/max was slow. Operationally, higher-effort tested configurations did
not guarantee better quality.

\subsection{Post-hoc Luna solo results}

\begin{center}
\begin{tabular}{@{}lrrrrr@{}}
\toprule
Exploratory arm & Checks & Mean & Median & Range & Timeouts \\
\midrule
Luna/xhigh solo & 110/120 & 750.9s & 780.5s & 621--816s & 0 \\
Luna/max solo & 109/120 & 779.4s & 768.3s & 682--909s & 0 \\
\bottomrule
\end{tabular}
\end{center}

Every run had exactly one observed Luna parent, no child session or Task call,
a normal exit, and only allowed file changes. Xhigh used 3.7\% less aggregate
time and scored one check higher than max, but was faster in only two of five
randomized blocks. Mean paired reduction was 3.1\%, descriptive 95\% interval
$-11.5$--17.7\%; mean savings were 28.5 seconds, interval $-86.2$--143.2.
There is no clear latency effect between these effort settings.

\begin{center}
\small
\begin{tabular}{@{}lrrr@{}}
\toprule
Non-concurrent baseline & vs Luna/xhigh & vs Luna/max & Quality \\
\midrule
Grok 4.5 solo & 85.4\% less time & 86.0\% less & 113 vs 110 / 109 \\
Sol/xhigh solo & 35.2\% less time & 37.5\% less & 112 vs 110 / 109 \\
\bottomrule
\end{tabular}
\end{center}

These historical gaps reinforce the operating choice but are not randomized
within-block effects because host and service time differed.

\section{Quality and Resource Analysis}

\subsection{Systematic misses}

No preregistered arm achieved 24/24, and all seven preregistered arms solved
comprehension perfectly:

\begin{center}
\begin{tabular}{@{}lrrr@{}}
\toprule
Arm & Implementation & Security & Comprehension \\
\midrule
Grok 4.5 & 35/40 & 38/40 & 40/40 \\
Grok 4.6 high & 35/40 & 37/40 & 40/40 \\
Grok 4.6 xhigh & 36/40 & 35/40 & 40/40 \\
Sol solo & 37/40 & 35/40 & 40/40 \\
Bounded Terra & 35/40 & 31/40 & 40/40 \\
Bounded Luna & 35/40 & 35/40 & 40/40 \\
Natural Luna & 35/40 & 35/40 & 40/40 \\
\bottomrule
\end{tabular}
\end{center}

Every Grok 4.5 run missed full invalid-concurrency validation; two also missed a
redirect parser edge. Every Sol run missed the redirect edge; three missed
invalid concurrency. Terra additionally missed profile-patch protections in two
runs. Public tests were not sufficient to certify completion. Contract-focused
tests remain mandatory, especially around nullability, parser layers,
authentication, redirects, cancellation, and cleanup races.

Neither exploratory Luna solo arm produced a perfect run. Both scored 35/40 on
implementation and 35/40 on security, systematically missing invalid-
concurrency and redirect checks. Luna/xhigh scored 40/40 on comprehension;
Luna/max scored 39/40 after one \texttt{dedupeKey} miss.

\subsection{Telemetry}

Across five Cursor runs, Grok 4.5 reported 135,151 input, 51,494 output, and
649,600 cache-read tokens. Grok 4.6 high reported 523,892, 143,195, and
1,432,192; xhigh reported 889,517, 197,714, and 2,333,568. Cursor emitted no
authoritative subscription charge or separate reasoning-token field.

OpenCode database telemetry and the prior nominal rate card gave:

\begin{center}
\small
\begin{tabular}{@{}lrrrrr@{}}
\toprule
Arm & Sessions & Input & Output & Reasoning & Proxy cost \\
\midrule
Sol solo & 5 & 212,008 & 37,775 & 79,682 & \$5.117 \\
Bounded Terra & 20 & 368,169 & 67,375 & 64,343 & \$4.720 \\
Bounded Luna & 20 & 691,841 & 97,013 & 303,225 & \$5.741 \\
Natural Luna & 21 & 890,564 & 102,395 & 403,134 & \$6.432 \\
Luna/xhigh solo* & 5 & 426,285 & 60,729 & 130,879 & \$1.896 \\
Luna/max solo* & 5 & 373,067 & 49,511 & 153,038 & \$1.891 \\
\bottomrule
\end{tabular}
\end{center}

The starred rows are from the post-hoc supplement. Their lower nominal proxy
reflects the Luna rate card, not faster completion or an actual invoice.

These are API-equivalent proxies, not OAuth subscription invoices, and Cursor
and OpenCode token semantics differ. Time, hidden quality, and rate-limit
behavior are more actionable than invented cross-vendor dollars.

\section{Protocol Deviations and Integrity}

All 35 outcomes exist, match the frozen order, share the fixture and prompt
digests, and match the stored grader outputs and current artifact hashes. A
manifest verified 163 raw
artifacts before reanalysis. Requested and observed models matched; Cursor made
zero Task calls; final changes were in scope. No exact grader filename or grader
directory appeared in a model stream.

Adversarial review found two bounded deviations:

\begin{itemize}
  \item Terra trial 2 temporarily edited \texttt{security/public.test.js}. The
  parent resumed the same worker in a fourth Task call and restored it before
  grading.
  \item Luna trial 4 resumed an implementation worker after it reported its
  configured step cap, weakening a strict 12-turn interpretation.
\end{itemize}

Both remain in intention-to-treat estimates and are excluded in clean-block
sensitivity estimates. Neither affects the Grok-versus-Sol contrast.

Isolation was policy-based, not OS-enforced. Workspaces were fresh and graders
were outside them, but Cursor used \texttt{--sandbox disabled}, processes
inherited the host environment, and harness credentials were present. No
observed path leakage is not proof of non-access. Future replication should use
containers or namespaces and broker credentials outside model-executed shells.

The post-hoc supplement retained ten outcomes plus streams, stderr, and grader
outputs. Its 43-file manifest verified before reanalysis. Every observed parent
was Luna with the requested xhigh or max variant; all runs were solo and in
scope. Its analyzer also verifies the parent 163-file manifest and records the
parent \texttt{runs.json} digest. Independent audit found no arithmetic or
current artifact-integrity error; methodological limits remain material. Source/
config hashes were captured after outcomes and are disclosed as such.

\section{Reconciling Old and New OpenCode Evidence}

\begin{center}
\begin{tabular}{@{}lrrrr@{}}
\toprule
Study & Sol mean & Terra mean & Reduction & Quality change \\
\midrule
August 12, OC 1.18.7 & 19.5m & 11.3m & 42.1\% & $+4$/120 corrected \\
August 14, OC 1.18.18 & 8.1m & 6.8m & 16.8\% & $-6$/120 \\
\bottomrule
\end{tabular}
\end{center}

The difference is evidence, not an inconvenience. Provider sampling, server
conditions, harness versions, and model behavior make agent-system effects
temporally unstable. One fixture estimates repeatability in one observation
window, not a permanent architecture constant.

The synthesis is narrower:

\begin{enumerate}
  \item parallel delegation can compress critical paths for independent work;
  \item the benefit can be modest or unreliable even on a favorable fixture;
  \item worker quality and parent correction determine whether speed survives;
  \item a faster solo harness/model can dominate the delegated architecture;
  \item architecture selection needs external baselines and periodic retesting.
\end{enumerate}

The August 15 solo supplement adds a same-harness descriptive check: Luna/xhigh
and Luna/max averaged 12.5 and 13.0 minutes versus the August 14 Sol mean of 8.1
minutes. Sol was not rerun concurrently, so this corroborates rather than
estimates a randomized Sol-versus-Luna effect.

\section{Adversarial Defense}

\subsection{One synthetic fixture cannot select a universal winner}

\textbf{Sustained.} This is a personal deployment decision. Multi-language,
SWE-bench-scale, coupled, interactive, and long-horizon replication is needed
for broad claims.

\subsection{The workload was engineered for subagents}

\textbf{Sustained, but conservative for the result.} Explicit independence
makes delegation easier. It cannot explain the solo arm's advantage.

\subsection{Harness and model are confounded}

\textbf{Sustained by design.} The user selects complete systems. A factorial
study would be needed to attribute the advantage to Grok versus Cursor's loop.

\subsection{Cursor had no turn cap}

\textbf{Sustained.} No equivalent was exposed. A common wall cap estimates
which deployable system manages work efficiently.

\subsection{Five blocks are too few}

\textbf{Sustained for population inference.} The Grok--Sol gap was large and
directionally consistent, but task-distribution coverage remains absent.

\subsection{Scores differ by one check}

\textbf{Sustained.} No intelligence superiority is claimed. Grok met quality
parity while taking far less time. Every system still requires verification.

\subsection{The timeout biases natural-Luna time}

\textbf{Sustained.} It is right-censored. The capped result is a lower bound on
Sol's advantage; four completed pairs still favor Sol by 51.9\%.

\subsection{Subscription economics and isolation are incomplete}

\textbf{Sustained.} Actual plan allowances should be monitored. A future study
needs OS isolation and credential brokering. This report does not manufacture
cross-vendor prices or claim proof of grader non-access.

\subsection{Luna solo was added after the parent result}

\textbf{Sustained.} The supplement is post-hoc relative to the parent decision
and cannot repair the omission retroactively. Its two-arm order was frozen
before those ten outcomes, so xhigh versus max is randomized within the
supplement. Grok and Sol comparisons remain non-concurrent and descriptive.

\section{Operating Policy}

\subsection{Default mode}

Use solo Cursor/Grok 4.5 high. Keep Task disabled when latency and attribution
matter. Give one coherent task, inspect the diff, and require tests.

\subsection{Verification gate}

\begin{enumerate}
  \item Run repository tests rather than trusting the completion summary.
  \item Add focused checks for contract boundaries omitted by public tests.
  \item Inspect changed-file scope, including tests and configuration.
  \item Review security parsing and cancellation paths independently.
  \item Reject ``done'' when final edits were not followed by rerun checks.
\end{enumerate}

\subsection{OpenCode mode}

Use Sol/xhigh solo with Task denied if OpenCode is needed. The post-hoc solo
screen does not support replacing Sol with Luna/xhigh or Luna/max. Do not expose
broad Luna agents by default. Preserve bounded Terra only as an explicit
experiment for low-risk independent packages with deterministic tests and
parent review; the current evidence does not support security-sensitive use.

\subsection{Rebenchmark triggers}

Repeat a frozen comparison after a major harness or model change, subscription
routing change, primary-language change, repository-scale change, or a shift
from independent packages to coupled refactors. Compare benchmark predictions
against actual completion and correction logs.

\section{Threats to Validity}

\subsection{Internal validity}

Sampling was nondeterministic. Harness, model, and effort were bundled. Two
bounded runs deviated. Five blocks cannot balance every time trend. Isolation
was not OS-enforced. The Luna supplement was requested after the parent result,
and its source hashes were captured post-outcome.

\subsection{Construct validity}

Hidden checks measure selected edges, not maintainability or complete security.
Wall time omits interactive steering. Nominal OpenCode cost is not billing.
Repeated fixture completion overrepresents deterministic tasks.

\subsection{External validity}

One JavaScript fixture, machine, account, and three-day window were studied.
Separability favors delegation. No large monorepository, UI task, migration,
dependency upgrade, or multi-hour task was tested. Cursor subagents were not
tested because the counterclaim concerned Cursor solo.

\subsection{Statistical conclusion validity}

Five blocks are imprecise for modest effects. Checks are clustered within runs.
Best-Cursor selection used a frozen quality-first rule but still evaluated
three variants on one fixture. The natural timeout is censored. Luna xhigh
versus max also has five blocks, and its historical comparisons are not paired.

\section{Reproducibility}

The private research archive contains the frozen manifest, protocol and pilot amendments,
post-outcome deviation ledger, isolated OpenCode config, resumable runner, 35
consolidated records, per-run raw streams and graders, a 163-file SHA-256
manifest, and standard-library-only analysis. The
\texttt{exploratory/luna-solo/} supplement contains its separate frozen two-arm
manifest, isolated config, ten outcomes, 43-file hash manifest, and analyzer.

\begin{lstlisting}[style=config]
cd path/to/reproducibility/cross-harness
node analyze.mjs
cd exploratory/luna-solo
node analyze.mjs
\end{lstlisting}

The analyzer verifies current manifest and fixture digests, recomputes package
scores from stored grader output, and checks count and order, prompt digests,
model identity, Cursor Task calls, grader exits, per-run/consolidated equality,
exact grader path appearances, final file scope, and raw artifact hashes. These
checks do not prove pre-outcome source immutability or OS-level non-access to
graders.

\section{Conclusion}

The original paper was right about mechanism but wrong to stop at the OpenCode
boundary. Bounded delegation can compress a critical path. That does not make
subagents the best workflow or a good default.

In the missing head-to-head comparison, solo Cursor/Grok 4.5 high was 77.5\%
faster than OpenCode Sol/xhigh solo, 73.0\% faster than bounded Terra, and more
accurate in aggregate than both. Luna/max subagents were substantially slower in
the tested delegated configurations. The requested follow-up also found Luna/xhigh and Luna/max solo slower and
lower-scoring than historical Sol and Grok; xhigh did not clearly beat max in
the randomized supplement. Higher-effort configurations did not improve the
decision outcome. The concern about subagent overhead is supported, and
switching the OpenCode solo owner from Sol to Luna is not supported either.

\resultbox{\textbf{Final thesis.} Optimize the complete coding system, not the
number of agents. Start with the fastest quality-eligible solo agent. Add
  delegation only after a task-specific hidden-graded comparison shows its
  integration overhead and quality risk are worth paying. On this fixture and
  observation window,
that means Cursor Agent with Grok 4.5 high, no subagents, and mandatory
verification.}

\section*{References}
\addcontentsline{toc}{section}{References}

\begin{enumerate}
  \item C. Jimenez et al., ``SWE-bench: Can Language Models Resolve Real-World
  GitHub Issues?'' ICLR 2024, arXiv:2310.06770.
  \item J. Yang et al., ``SWE-agent: Agent-Computer Interfaces Enable Automated
  Software Engineering,'' arXiv:2405.15793.
  \item C. S. Xia et al., ``Agentless: Demystifying LLM-based Software
  Engineering Agents,'' arXiv:2407.01489.
  \item Cursor, Headless CLI, \url{https://cursor.com/docs/cli/headless}.
  \item Cursor, Output format,
  \url{https://cursor.com/docs/cli/reference/output-format}.
  \item Cursor, Subagents, \url{https://cursor.com/docs/subagents}.
  \item OpenCode, Agents, \url{https://opencode.ai/docs/agents/}.
  \item OpenCode, Models, \url{https://opencode.ai/docs/models/}.
  \item Historical local report,
  \texttt{opencode-subagent-performance-audit-2026-08-12.md}.
\end{enumerate}

\end{document}
